\clearpage
\section{Stable graphs and copyable payloads}
\label{sec:copying-values}
Definitions~10.2 and 10.3 serve different purposes. Copyability describes
what Figure~17 may duplicate; stability restricts the graph used by the
value-sensitive rules. We now formalise both for the paper's full example
operation alphabet. Neither the proposition terms nor their Figure~16
translation has changed. In particular, extracting a copyable payload below
is an exact factorisation of the original translated graph, not a replacement
of that graph by an equivalent program.

These are structural predicates and proofs. Their use in copying, store
transitions and the operational stability/refocusing invariants remains
unfinished. The existing correspondence audit of the earlier linear
substitution rule remains open.

\subsection{The complete stability predicate}
\termdef{copying-stable-graph}
A stable graph has a single term input; every edge is an instance or a
passive operation; every external output is an atom or a thunk; and every
vertex is reachable from that input. Boxes are not stable edges: a lambda
head may expose a thunk output to which a box is attached outside the stable
fragment. The definition retains all three conditions of Definition~10.3,
rather than replacing reachability with a language-label approximation.

The primitive predicate accepts an instance and accepts an operation exactly
when its declared kind is passive. Its lifted label predicate requires an
actual primitive label, excluding boxes even if their contents happen to
contain passive operations. The output predicate accepts atoms and thunks
and rejects terms. A vertex step follows one edge from any source to any
target. Reachability is existence of an exact finite trace of such steps,
including the zero-step trace. The stability structure then records the
label, output, and all-vertex reachability requirements explicitly.

\lean{Copying/Stability}{12}{15}{Stability.StablePrimitive}

\lean{Copying/Stability}{18}{19}{Stability.StableLabel}

\lean{Copying/Stability}{22}{24}{Stability.StableOutput}

\lean{Copying/Stability}{27}{28}{Stability.VertexStep}

\lean{Copying/Stability}{30}{31}{Stability.VertexReach}

\lean{Copying/Stability}{34}{39}{Stability.Stable}

Every permitted primitive has one term input. In the example alphabet,
its outputs are non-term: a constant has no outputs, lambda has a thunk
output, and an instance has an atom output. Case analysis on the primitive
and operation constructors proves these signature facts. Mapping labels
over source and target incidence lists transfers them to actual vertices:
a source vertex of a permitted edge has term type, whereas a target vertex
has a permitted non-term type.

\lean{Copying/Stability}{41}{43}{Stability.primitive_input}

\lean{Copying/Stability}{45}{50}{Stability.primitive_outputs}

\lean{Copying/Stability}{52}{58}{Stability.source_type}

\lean{Copying/Stability}{60}{66}{Stability.target_type}

An internal vertex must occur in some source and some target, by the
graph's two coverage laws. The preceding typing facts would therefore
make it both term and non-term. Thus no internal vertex exists. Each edge
has a nonempty singleton source list. Its source cannot be an output by
boundary discipline or an internal vertex by the preceding contradiction,
so it is the unique external input. Source-incidence uniqueness now makes
any two edges equal. Coverage at the input also supplies an edge, proving
existence as well as uniqueness. These conclusions use the concrete example
signatures; they are not asserted for an arbitrary operation alphabet.

\lean{Copying/Stability}{69}{75}{Stability.no_internal}

\lean{Copying/Stability}{78}{92}{Stability.input_source}

\lean{Copying/Stability}{95}{99}{Stability.edge_unique}

\lean{Copying/Stability}{100}{103}{Stability.exists_unique_edge}

For this alphabet, the output and reachability conditions follow from
the permitted labels. Each output has a producing edge by target coverage,
so its type is permitted. For reachability, the input reaches itself in
zero steps, internal vertices are impossible, and each output is reached
in one step through its producer, which consumes the input. These proofs
construct the full stability record rather than assuming the extra fields.
Finally, an adjacency would share a vertex between a target and a source,
again forcing a non-term vertex to be term. Hence there is no outer edge
adjacency in a stable graph of this alphabet.

\lean{Copying/Stability}{106}{120}{Stability.of_labels}

\lean{Copying/Stability}{123}{127}{Stability.no_adjacency}

Constant, lambda-head, and instance corollas satisfy the label restriction;
the construction just proved supplies their output and reachability fields
at every box allowance. Stability also implies global focus freedom and
hole freedom. Every outer label is a primitive instance or passive operation,
so it cannot equal either forbidden constructor. No box case is hidden:
the stability label predicate excludes boxes outright.

\lean{Copying/Stability}{130}{135}{Stability.constant}

\lean{Copying/Stability}{137}{142}{Stability.lambda}

\lean{Copying/Stability}{144}{149}{Stability.instanceEdge}

\lean{Copying/Stability}{151}{156}{Stability.focusFree}

\lean{Copying/Stability}{158}{163}{Stability.holeFree}

Recursive label equivalence preserves the stable-label predicate. At depth
zero it is literal equality. At positive depth, related primitive labels
are equal; primitive/box pairs cannot be related; and neither box satisfies
the stable-label predicate. To transport stability through a recursive graph
isomorphism, pull each target edge back through the edge bijection, transfer
its stable label, and invoke the proved construction from labels. Applying
the same argument to the inverse gives equivalence of stability.

\lean{Copying/Stability}{166}{177}{Stability.label_equiv}

\lean{Copying/Stability}{179}{183}{Stability.iso}

\lean{Copying/Stability}{185}{186}{Stability.iso_iff}

\subsection{Copyability specialised to the full example alphabet}
\termdef{copying-copyable-payload}
Definition~10.2 permits an instance, or an operation followed by the parallel
box forest attached to its thunk outputs. In $\mathbb O^{\mathrm{ex}}$, only
lambda has a thunk output, and it has exactly one. Thus lambda has exactly
one attached box; constants, application, unary operations and binary
operations have no attached boxes. Specialising the diagram gives the
constructors below without removing any case of the example language.

A lambda payload is serial assembly of the lambda corolla with the box
of its body, before adding the translation's outer distributors. The box
has one bound variable and the original ordered free variable/atom interface.
Copyability has constructors for instance, constant, application, each unary
operation, each binary operation, and this lambda payload. A final constructor
closes these shapes under recursive graph isomorphism, including independent
renaming inside boxes. The predicate retains the interface counts of each
shape. Box bodies are not assumed stable: they may contain computations.
Copyability is a shape predicate, not an assertion of termination.

\lean{Copying/Copyability}{10}{15}{Copyability.lambdaPayload}

\lean{Copying/Copyability}{16}{26}{Copyability.Copyable}

The outer-label predicate allows only instance, operation, and box labels.
It is invariant under recursive label equivalence by the same primitive/box
case analysis. Induction on a copyability derivation proves that every outer
edge has an allowed label. Primitive cases inspect the corolla; the lambda
case splits between its lambda and box edges; the renaming case pulls back
through the edge bijection and uses label invariance. Consequently no outer
sharing, store, focus or hole edge is part of a copyable payload. This is
an outer-level statement; arbitrary box bodies require their own invariants.
Copyability is preserved and reflected by recursive isomorphism, using its
renaming constructor in both directions.

\lean{Copying/Copyability}{29}{35}{Copyability.OuterLabel}

\lean{Copying/Copyability}{37}{48}{Copyability.outer_equiv}

\lean{Copying/Copyability}{51}{63}{Copyability.outer_labels}

\lean{Copying/Copyability}{65}{70}{Copyability.iso_iff}

The literal abstraction translation factors as the lambda payload followed
by $D^{\star}_{k,1}\otimes D^{\diamond}_{h,1}$. The equation is definitional:
it only names an existing part of the assembled graph. In particular,
none of the distributors is deleted. If the body is a language graph, the
payload is a language graph too, by the box-language lemma, the permitted
lambda primitive, and serial closure. This supplies the stronger global
invariants when the body comes from the term translation.

\lean{Copying/Copyability}{71}{72}{Copyability.abstraction_factor}

\lean{Copying/Copyability}{75}{77}{Copyability.lambdaPayload_language}

\subsection{Distinguishing the predicates and using the original program}

Every unary operation corolla is copyable, including reference creation
and dereference, but it is not stable: its term output contradicts the
stable-output field. A lambda payload is also not stable, even though its
head is passive. Serial assembly introduces an internal thunk vertex between
lambda and box, contradicting the proved absence of internal vertices in
a stable graph. These negative results prevent treating ``copyable'' and
``stable'' as interchangeable approximations to values.

For every scoped term supplied as a lambda body, its translated payload
is copyable by construction. The full translation's language theorem and
payload language preservation prove global focus and hole freedom. Thus
arbitrary computations may remain inside the box while the fragment used
by the copy rule has the required structural form.

\lean{Copying/ValueExamples}{9}{14}{ValueExamples.unary_copyable_not_stable}

\lean{Copying/ValueExamples}{17}{21}{ValueExamples.lambdaPayload_not_stable}

\lean{Copying/ValueExamples}{24}{32}{ValueExamples.translated_lambda_payload}

For every ambient $(k,h)$, the translation of the unchanged right-hand
program $\lambda f.(f();1)$ factors into a copyable payload and its two
single-use distributors. The witness is the lambda payload of the translated
body $f();1$, using exactly the existing sequencing sugar. The factorisation
equality is reflexivity after unfolding the construction; no contextual
or semantic equivalence is assumed. The graph has allowance two because
sequencing introduces the second abstraction. Finally, the unit corolla
is stable and has a unique edge by the general stability theorems.


\lean{Copying/ValueExamples}{33}{38}{ValueExamples.constantResult_copyable_factor}

\lean{Copying/ValueExamples}{41}{44}{ValueExamples.unit_stable_unique}

All three axiom audits and the full build pass with only standard Lean
axioms. The definitions and structural properties of stable and copyable
graphs are completed. Required depth and operational preservation laws,
copying/store transitions, translation/context laws, and template robustness
remain in Appendix~\ref{app:remaining-work}. Proposition~12.2 is not yet proved.
